%% kapa: exam=2 round=27 year=2016 period=1 subject=appraisal_practice subject_ko=감정평가실무 seq=3 number=3 points=20
%% source: Q-Net 감정평가사 기출문제 (cst00309) artlSeq=5185401 file=extracted/1008/5185401/제27회 감정평가사 2차시험 기출문제/제27회 감정평가사 2차 시험문제지/(1교시 시험문제지)감정평가실무.pdf
%% conversion: kordoc 4.24.51 -> tools/gen2.py (auto); review=proofread-gpt-sol; flags=none
\begin{longquestion}{【문제 3】}
감정평가사 甲은 국산 사출기 20대를 보유하고, 플라스틱 제품을 생산하여, 수출 중인 사업체(K사) 전체에 대한 적정한 시장가치의 산정을 의뢰받았다. 토지, 건물 및 구축물, 영업권 등의 무형자산에 대한 가치까지 산정한 후, 최종적으로 사업체의 주 생산설비인 국산 사출기 20대에 대하여 관련 규칙 및 기준에 의거하여 평가하고자 한다. 제시된 자료를 참조하여 평가방법을 결정하고, 다음 물음에 답하시오. (20점) \par
\medskip\noindent 물음 1) ‘감정평가에 관한 규칙’ 별지서식인 “감정평가액의 산출근거 및 결정 의견”을 최대한 활용하여 제1라인의 적정가격을 제시하시오. (10점) \par
\medskip\noindent 물음 2) ‘감정평가에 관한 규칙’ 별지서식인 “감정평가액의 산출근거 및 결정 의견”을 최대한 활용하여 제2라인의 적정가격을 제시하시오. (10점) \par
<자료 1> 기본사항 \par
1. 의뢰인 및 사업체명 : 주식회사 K \par
2. 기준시점 : 2016년 7월 1일 \par
3. 생산라인구성 : 제1라인과 제2라인으로 구성되어 있으며, 각 생산라인에 10대씩 설치되어 있으나, 제1라인의 사출기는 생산효율이 높지 아니하고 전용불가능한 과잉유휴설비로 전체를 철거하여 매각할 예정임 \par
<자료 2> 기계에 관한 사항 \par
1. 제1라인 : 2006년 7월 1일 10대 설치가동, 유지보수상태 보통이하 \par
2. 제2라인 : 2011년 7월 1일 10대 설치가동, 유지보수상태 양호 \par
<자료 3> 라인별 취득가격 및 유지보수비 등 \par
1. 라인별 취득가격 \par
{\small\centering
\resizebox{\ifdim\width>\linewidth\linewidth\else\width\fi}{!}{%
\begin{tabular}{|c|c|c|}
\hline
\makecell{구 분} & \makecell{제1라인 단위당 취득가격(원)} & \makecell{제2라인 단위당 취득가격(원)} \\
\hline
\makecell{본체} & \makecell{50,000,000} & \makecell{80,000,000} \\
\hline
\makecell{부대설비} & \makecell{20,000,000} & \makecell{30,000,000} \\
\hline
\makecell{설치비} & \makecell{5,000,000} & \makecell{5,000,000} \\
\hline
\makecell{시험운전비} & \makecell{5,000,000} & \makecell{5,000,000} \\
\hline
\makecell{부가가치세} & \makecell{8,000,000} & \makecell{12,000,000} \\
\hline
\end{tabular}}\par}

2. 제1라인의 경우, 설치이후 현재까지 단위당 유지보수 등을 위한 수익적 지출에 20,000,000원, 자본적 지출에 20,000,000원이 각각 소요됨 \par
3. 제2라인의 경우, 설치이후 현재까지 단위당 유지보수 등을 위한 수익적 지출에 10,000,000원, 자본적 지출에 10,000,000원이 각각 소요됨 \par
<자료 4> 내용연수 및 잔가율 등 \par
1. 국내생산의 사출기는 물리적 내용연수가 12년, 경제적 내용연수는 10년 정도인 것으로 조사됨 \par
2. 본 기계의 잔가율은 통상 10 \%로 조사되고, 감가수정은 관련 법령에서 제시한 원칙적 방법에 따를 예정임 \par
3. 물가변동에 따른 기계가격 보정지수 : 취득가격에만 적용 \par
1) 제1라인은 기준시점까지 10 \% 상승 \par
2) 제2라인은 기준시점까지 변동사항이 없음 \par
<자료 5> 기타자료 \par
1. 제1라인의 유사사양 사출기는 생산효율의 저감으로, 해체 및 포장된 상태에서 동남아 등지에 기계를 수출하는 업자에게 매각가능하며, 단위당 매각가능가격은 잔존가치와 유사한 것으로 조사됨 \par
2. 제1라인의 해체 및 철거와 조립 및 포장 운반 등에 소요되는 단위당 관련 비용은 아래와 같이 조사됨 \par
1) 해체비 : 1,000,000원 \par
2) 철거비 : 1,000,000원 \par
3) 운반비 등 : 1,000,000원 \par
4) 설치비 : 5,000,000원 \par
\end{longquestion}
